Im Zylinder eines Heissluftmotors befinden sich zunächst \VEiO\ Luft bei einer Temperatur von \thEiO\ und einem Druck von \pEiO. Durch den Überdruck wird das Volumen isotherm auf \VZwO\ vergrössert. In der neuen Position wird die Luft isochor auf \thDrO\ abgekühlt und anschliessend isotherm wieder auf das Anfangsvolumen verdichtet.
\begin{abcliste}
\abc Berechnen Sie die Stoffmenge und die Masse der Luft im Zylinder (molare Masse von Luft: \MLuO).
\abc Erstellen Sie ein $p(V)$-Diagramm, in dem Sie die vier beschriebenen Zustände einzeichnen.
\end{abcliste}
